\chapter{Motivating Study: Intent-Conditioned Single-Path Selection}
\label{ch:single_path}

This chapter investigates the first and most basic control decision a SCION endpoint faces: given a
running application's intent and a changing network, \emph{which} of the candidate paths to a destination
should carry a flow? We answer it in five steps. \Cref{sec:sp:framework} builds the selection framework --
the simulation environment, the decision problem, and a permutation-equivariant scoring architecture that
copes with a variable path set. \Cref{sec:p1design:film} adds the chapter's central contribution, FiLM-based
intent conditioning, and argues why modulation is required where naive concatenation fails.
\Cref{sec:sp:impl} describes the implementation and the baselines. \Cref{sec:sp:eval} evaluates the
conditioned selector on held-out hours. Finally, \cref{sec:p1eval:ceiling} identifies
the single-path performance ceiling that motivates the multipath system of \cref{ch:multipath}.

\section{Selection Framework: Environment, Decision, and Scoring}
\label{sec:sp:framework}

We start from the DQN path selector of Bourak~et~al.~\cite{Bourak2025} and identify three limitations we
set out to remove (\cref{sec:p1design:strawman}). We then describe the improved simulation environment
(\cref{sec:p1design:env}), the state, action, and reward that define the decision problem
(\cref{sec:p1design:mdp}), and the permutation-equivariant path-scoring architecture that handles a
variable path set (\cref{sec:p1design:scoring}).

\subsection{Starting Point and Its Limitations}
\label{sec:p1design:strawman}

Bourak~et~al.~\cite{Bourak2025} frame SCION path selection as a reinforcement-learning problem and
train a DQN to pick a path from per-path measurements of latency, loss, and bandwidth, rewarding a
weighted combination of goodput and trust while charging for probing. They report oracle-level
selection quality at a fraction of the probing cost of measuring every path. This validates the RL
framing and gives us a reward structure to build on, but three design choices limit it.

\begin{description}
  \item[Fixed action space.] A vanilla DQN has one output per action, so the number of paths must be
  fixed and known. Real source--destination pairs have different, time-varying path counts; padding to
  a maximum wastes capacity and still caps the topology size.
  \item[Single baked-in objective.] The reward weights are constants chosen before training. The
  resulting policy embodies \emph{one} trade-off between goodput and trust; a different application
  intent demands a different, separately trained model.
  \item[Environment fidelity.] Capturing when path selection actually matters requires paths that
  genuinely share bottlenecks and conditions that vary over realistic timescales, so that a good choice
  at one hour is a poor choice at another.
\end{description}

Our design removes all three: a scoring architecture for the first, FiLM conditioning for the second,
and an improved environment for the third. A strawman for the second limitation -- appending the reward
weights to the state of the existing DQN -- is instructive precisely because of \emph{where} it fails, as
we show in \cref{sec:p1design:film}; that failure motivates the modulation-based design.

\gap{the third limitation is asserted on both sides and measured on neither: the draft claims Bourak
et~al.'s environment lacks the fidelity to make selection matter, and that ours has it, without a number
supporting either. One statistic would settle it and belongs in \cref{sec:p1design:env} -- the spread of
achievable goodput across a pair's candidate paths within a decision context, which is exactly the
quantity a selector has to exploit. If that spread is wide and varies by hour, path choice demonstrably
matters here; if it is narrow, every result in this chapter is a small effect. The run's characterization
files already hold the supporting evidence for a short environment-realism table: p90 link utilization
0.76 with 7.8\% of links oversubscribed, 6.8\% of pair-hours with no usable bandwidth on any candidate,
and 29.95 candidate paths per pair averaged over 8010~routable pairs.}

\subsection{Improved SCION Simulation Environment}
\label{sec:p1design:env}

Training and evaluating a path selector requires an environment in which the ``right'' path changes
with conditions. We construct one in several stages.

\paragraph{Topology.} AS-level topologies are generated with BRITE and converted into SCION structures
-- ISDs, core and non-core ASes, and inter/intra-ISD links -- so that the graph has realistic degree
and locality structure rather than being hand-drawn.
\gap{describe the conversion honestly and concretely. BRITE is run with a Barab\'{a}si--Albert model on
50~nodes with every generated link at exactly \SI{10}{\giga\bit\per\second}; the 90~ASes and 528~links
of \cref{sec:p1eval:setup} arise because the converter then \emph{adds} synthetic structure with
hard-coded capacities -- a full intra-ISD core mesh at \SI{100}{\giga\bit\per\second}, an inter-ISD core
ring at \SI{50}{\giga\bit\per\second}, multi-parent and cross-connect links at
\SIrange{1}{2}{\giga\bit\per\second}, and peer links drawn uniformly from
\SIrange{5}{10}{\giga\bit\per\second}. ISDs are assigned by $k$-means on node coordinates and core ASes
by degree rank. Link latency is \emph{not} BRITE's delay column but $\max(0.5,\,0.02\,d)$ ms in the node
coordinate plane. State this, because the capacity heterogeneity the agent exploits is largely an
artifact of the converter rather than of BRITE, and a reader who assumes BRITE-generated bandwidths will
misread the results. The three pipeline figures
\texttt{topology/step\{1,2,3\}\_*.png} in the run directory illustrate exactly this
BRITE~$\rightarrow$~SCION~$\rightarrow$~peering progression and are already publication quality; include
one of them here.}

\paragraph{Control plane.} A beaconing simulation propagates path-segment construction beacons
top-down within each ISD and across the core mesh between ISDs, yielding the up/core/down segments that
compose into end-to-end paths. Peer links are excluded. The result is, for each source--destination
pair, a realistic set of candidate paths of varying length and quality -- the agent's action set.

\paragraph{Traffic and link dynamics.} We drive the network with 28~days of hourly demand: foreground
flows for every routable pair plus randomized background traffic, modulated by diurnal and weekly
factors. Crucially, load is aggregated \emph{per link}, so paths that share a bottleneck link contend
for its capacity. This is what makes selection non-trivial: the utilization and effective bandwidth of
a path depend on what every other flow is doing at that hour, and the best path at 3\,a.m. is often
congested at 8\,p.m.

\paragraph{Link and path model.} A link's utilization is its aggregate offered load over its capacity,
capped at 1.5; its latency is a base propagation delay scaled by a queueing multiplier that grows with
utilization and saturates at $2\times$; and its loss rate is a piecewise function of utilization alone,
zero below 0.8 and capped at 0.2. A path's latency is the sum over its links, its loss the complement of
the product of per-link survival probabilities, its utilization the maximum over its links, and its
\emph{available bandwidth} the minimum residual capacity $\max(0,\,C_e - \mathrm{load}_e)$ across its
links.

\gap{state the two consequences of this model that bound how the chapter's results should be read.
First, the model is \emph{open-loop}: the agent's own selection never adds load to the links it
traverses, so what the reward calls goodput is the spare capacity at the bottleneck rather than a rate
the flow would actually achieve, and there is no per-flow fair share or feedback. Second, because
utilization is capped at 1.5 rather than 1.0, links can be oversubscribed -- which is what produces the
$>1.0$ x-axis values in \cref{fig:p1eval:ceiling}. Both are defensible modeling choices for a selection
study, but they must be declared, not discovered by the reader at the ceiling plot.}

\paragraph{Probing model.} Fresh metrics are not free. The environment charges a latency probe
\SI{10}{\milli\second} plus \SI{0.5}{\milli\second} per hop, and a full bandwidth probe
\SI{100}{\milli\second} plus \SI{20}{\milli\second} per hop; it credits the agent with having probed
only the path(s) it actually inspects. Because both costs grow with hop count, probing is cheaper on
short paths -- a fact \cref{sec:p1eval:overhead} returns to. The reward charges the \emph{mean} cost per
probe issued in a step, normalized by a \SI{500}{\milli\second} reference. This turns ``measure less''
into a real objective rather than an afterthought, and lets us compare the learned selector's probing
footprint against baselines that probe exhaustively.

\gap{the per-probe normalization deserves a sentence of justification, because it is consequential and
non-obvious: charging the mean rather than the total cost means a method that issues 30~cheap probes is
penalized \emph{less} than one issuing a single expensive probe. That inverts the intuition the reward is
meant to encode and it directly affects the reward comparison in \cref{sec:p1eval:overhead}. Either
justify it (it measures cost per unit of information gained) or report the total-cost variant alongside;
the environment already supports switching the normalization off.}

\subsection{Decision Problem: State, Action, Reward}
\label{sec:p1design:mdp}

\paragraph{State.} The agent observes a \emph{global} context and a \emph{per-path} feature matrix. The
global context summarizes time and aggregate conditions -- normalized day-of-week and hour-of-day, mean
path utilization, mean trust, and the fraction of paths whose utilization exceeds 0.7 -- plus a compact
encoding of the source and destination. The three aggregate terms are computed over the current flow's
own candidate path set, not over the whole network. Each candidate path contributes a seven-dimensional
feature row: latency normalized by \SI{100}{\milli\second}, loss rate, hop count over 20, the path's
bandwidth relative to the best candidate at that hour, utilization, a static-bandwidth term, and a trust
score. Because the number of rows $N$ varies, the state is naturally a (global, path-matrix) pair rather
than a flat vector.

\paragraph{Action.} The action is the choice of one path, i.e.\ an index into the current candidate set.
Because $N$ varies, the action is realized as an $\arg\max$ over per-path scores (\cref{sec:p1design:scoring})
with invalid entries masked, rather than as a fixed categorical head.

\paragraph{Reward.} We adopt and generalize the goodput--trust reward of
Bourak~et~al.~\cite{Bourak2025}. For the chosen path, let $\mathrm{goodput}\in[0,1]$ be the achieved
bandwidth normalized by the \emph{best available bandwidth among that pair's candidate paths at that
hour}, and define a trust term from loss and delay,
\begin{equation}
\mathrm{trust} = \max\!\big(0,\; 1 - (w_3\,\mathrm{loss} + w_4\,\mathrm{delay})\big),
\label{eq:p1trust}
\end{equation}
with $\mathrm{delay}$ the latency capped at \SI{100}{\milli\second} and normalized. The base reward
trades goodput against trust and is shifted to $[-1,1]$, and a probing penalty is subtracted:
\begin{equation}
r = \underbrace{2\big(w_1\,\mathrm{goodput} + w_2\,\mathrm{trust}\big) - 1}_{\text{base reward}}
    \;-\; w_{\text{probe}}\cdot \mathrm{probe\_penalty},
\label{eq:p1reward}
\end{equation}
clamped to $[-1,1]$, where $\mathrm{probe\_penalty}$ is the probing cost normalized to $[0,1]$. We take
$w_1 + w_2 = 1$ throughout, which is what places the base reward in $[-1,1]$ before clamping. The
five weights $\mathbf{w} = (w_1, w_2, w_3, w_4, w_{\text{probe}})$ are exactly the \emph{application
intent}: default balanced values are $(0.7, 0.3, 0.5, 0.5, 0.05)$, but a throughput-hungry intent
pushes $w_1$ toward~1, a loss-averse intent raises $w_2$ and $w_3$, a latency-sensitive intent raises
$w_4$, and a probe-frugal intent raises $w_{\text{probe}}$. The design goal of \cref{sec:p1design:film}
is a \emph{single} policy that behaves correctly for \emph{any} setting of $\mathbf{w}$.

\gap{two properties of this reward shape every result in the chapter and should be stated here rather
than left implicit. First, because goodput is normalized against the \emph{best candidate path at the
same hour}, it is a relative score, not an absolute rate: a policy that always picks the best available
path scores near~1 whether that path carries \SI{9}{\giga\bit\per\second} or
\SI{2}{\giga\bit\per\second}. This is why the FiLM selector's reward stays essentially flat across the
congestion range in \cref{sec:p1eval:ceiling} while its goodput falls by a quarter -- the reward
measures selection quality, and the ceiling is invisible to it. Say so explicitly, because a reader who
takes the reward as an absolute performance measure will find \cref{sec:p1eval:ceiling} contradictory.
Second, latency enters the objective only \emph{inside} the trust term, via $w_4$; there is no
standalone latency weight. A ``latency-sensitive'' intent is therefore expressed by raising $w_2$ and
$w_4$ together, which is not obvious and is worth one sentence.}

\subsection{Handling a Variable Path Set: The Scoring Architecture}
\label{sec:p1design:scoring}

To make the policy independent of the path count, we score each path with a \emph{shared} network and
select by comparison, following the permutation-equivariant design of \cref{sec:background:variable}.
The global context is broadcast and concatenated to every path's feature row, and a shared multilayer
perceptron maps each $(\text{global}, \text{path}_i)$ pair to a scalar $Q$-value. Selection is
$\arg\max_i Q_i$; absent or invalid paths are masked to $-10^9$ so they can never be chosen and do not
distort training targets.

We use a \emph{dueling} scoring head (\cref{sec:background:rl}): a value stream reads the global context
alone to produce $V(s)$, an advantage stream reads the per-path features to produce $A(s, i)$, and the
two combine as $Q_i = V(s) + \big(A(s,i) - \operatorname{mean}_{j\in\text{valid}} A(s,j)\big)$ with the
mean taken only over valid paths. Because many paths are often comparably good, separating ``how good is
this situation'' from ``which path is marginally better'' stabilizes learning. Double DQN targets and
prioritized replay (\cref{sec:background:rl}) further stabilize and accelerate training. The scoring
design also delivers the dynamic-action property directly: the same weights serve a pair with three
paths and a pair with thirty, and are invariant to how the paths happen to be ordered.

\section{Intent Conditioning via Feature-Wise Linear Modulation}
\label{sec:p1design:film}

The central contribution of this chapter is making one scoring policy serve every intent $\mathbf{w}$. We
condition the network on $\mathbf{w}$, but the \emph{mechanism} of conditioning is what makes or breaks
it.

\subsection{Where conditioning cancels}
\label{sec:p1design:concat}
The obvious approach is to append $\mathbf{w}$ to the state. Whether that works depends entirely on
\emph{which} part of the state it is appended to, and the reason is a one-line argument about the
comparison that defines the action.

\begin{proposition}[Common-mode cancellation]
\label{prop:p1cancel}
Let a scoring policy assign path~$i$ the score
$Q_i(s, \mathbf{w}) = f(s, f_i) + c(s, \mathbf{w})$, where $f_i$ is path~$i$'s feature row and the
conditioning enters only through the term $c$, which does not depend on~$i$. Then
$\arg\max_i Q_i(s,\mathbf{w})$ is independent of $\mathbf{w}$: the policy consumes the intent but its
action is invariant to it.
\end{proposition}

\noindent The proof is immediate --- $c(s,\mathbf{w})$ is a constant offset across the compared
alternatives, so it cancels in every pairwise comparison $Q_i - Q_j$ and cannot change which index
attains the maximum. The point is that this is not a statement about optimization difficulty or sample
efficiency: no amount of training can overcome it, because the intent-dependent term is not in the image
of the decision rule at all.

The dueling scoring head of \cref{sec:p1design:scoring} instantiates exactly this hypothesis when the
intent is routed into the value stream. There $Q_i = V(s, \mathbf{w}) + \big(A(s,i) -
\operatorname{mean}_j A(s,j)\big)$, so $V$ plays the role of $c$: conditioning the value stream is
provably inert. \Cref{sec:p1eval:ablation} confirms this empirically, and we verify it directly at the
network level -- holding path features fixed and varying only $\mathbf{w}$ leaves every pairwise
$Q$-difference constant to within $10^{-5}$.

Two things follow, and it is worth being precise about both because the distinction is easy to lose.
First, the failure is a property of \emph{where} the signal enters, not of concatenation as a technique:
if $\mathbf{w}$ is concatenated onto each path's feature row and passed through a nonlinear scorer, the
resulting $Q_i(s, \mathbf{w}, f_i)$ is not of the form assumed in \cref{prop:p1cancel}, the argument does
not apply, and such a policy can re-rank. \Cref{sec:p1eval:ablation} measures one. Second, what
distinguishes the design of the next section is therefore not that it is the only mechanism that
\emph{can} re-rank, but that it is the only one that re-ranks \emph{by construction}: the intent's
effect on the relative ordering is structural rather than something the network must discover.

\subsection{Weight-FiLM dueling head}
The fix is to let the intent modulate the \emph{per-path} features multiplicatively, so that it changes
the relative ranking of paths rather than a shared baseline. A small conditioning network maps the intent
vector to feature-wise scale and shift parameters,
\begin{equation}
(\gamma, \beta) = \mathrm{FiLM}_\theta(\mathbf{w}), \qquad
\tilde{f}_i = (1 + \gamma) \odot f_i + \beta,
\label{eq:p1film}
\end{equation}
where $f_i$ is path~$i$'s feature row and $\tilde{f}_i$ the modulated features fed to the advantage
stream. A latency-sensitive intent can learn a $\gamma$ that amplifies the latency feature, so paths
differ more sharply in their latency contribution and a low-latency path wins; a throughput-hungry intent
can instead amplify bandwidth. Because the modulation is per-feature and multiplicative, it acts
\emph{differently} on paths with different feature values and therefore reorders them by construction --
exactly what \cref{prop:p1cancel} rules out for value-stream conditioning. Concretely, $\gamma$ and
$\beta$ are seven-dimensional (one scale and one shift per path feature), produced by a two-layer
conditioning network from the encoded intent, and shared across the paths of a decision -- it is the
paths' differing feature \emph{values}, not differing modulation, that makes the effect asymmetric. The
value stream continues to read only the topology/time context, keeping the state-value estimate
intent-agnostic while the advantage -- the part that decides \emph{which} path -- is intent-modulated.

\subsection{Training across intents}
To learn a policy that spans the intent space we train across six \emph{reward profiles} -- named
weight vectors spanning a throughput extreme, a loss-averse extreme, a latency-averse extreme, a balanced
midpoint, and two probe-frugality settings (\cref{tab:app:profiles}). Each training episode is assigned a
profile, sets the environment's reward weights to it, and feeds the same weights to the FiLM head, so the
agent simultaneously experiences the reward for an intent and the encoding of that intent. Profiles are
scheduled in a stratified fashion -- each profile repeated to fill the episode budget, then shuffled --
so every intent receives an equal number of episodes. The intent is presented to the network in a
rescaled encoding rather than as raw weights, which keeps the conditioning inputs on a comparable scale
to the path features.

\gap{the design intent is that presenting an unseen $\mathbf{w}$ at inference -- in particular an
interpolation between two trained profiles -- yields correspondingly interpolated behavior without
retraining. This is the \emph{zero-shot generalization} property claimed in \cref{sec:intro:contributions}
and required by \cref{sec:problem:properties}, and it is currently untested: every intent reported in
\cref{sec:sp:eval} is one of the six profiles the agent trained on. Either add the experiment
(\cref{sec:p1eval:zeroshot}) or drop the claim from the abstract, the contributions, and
\cref{sec:analysis:generalization}.}

\subsection{Relation to Mortise}
This is the path-selection instance of Mortise's principle~\cite{Shen2026}. Where Mortise maps an
application's revealed QoS preference to congestion-control \emph{parameters}, our FiLM head maps an
explicit intent vector to a \emph{modulation} of path features that determines which path is chosen. In
both, intent is a first-class input that reshapes a network mechanism, not a constant compiled into it.
\Cref{ch:multipath} pushes the same principle one step further, letting intent shape not only the network's
choice but the application's own output.

\section{Implementation}
\label{sec:sp:impl}

This section describes how the design is realized: the agent family and their architectures
(\cref{sec:p1impl:agents}), the end-to-end evaluation pipeline that turns a topology into trained models
and comparable results (\cref{sec:p1impl:pipeline}), and the baselines and probing accounting against
which the learned selectors are measured (\cref{sec:p1impl:baselines}). The implementation is in Python
(PyTorch for the agents) and is organized so that each pipeline stage passes data to the next through
timestamped run directories.

\subsection{Agent Family}
\label{sec:p1impl:agents}

The codebase implements a family of agents of increasing capability, from which the ablation ladder of
\cref{sec:p1eval:ablation} is drawn.

\begin{description}
  \item[Flat DQN (simple / enhanced).] A conventional fixed-action DQN over the 5-dimensional global
  state. The \emph{simple} variant is a small multilayer perceptron with $\varepsilon$-greedy
  exploration and a periodic target network; the \emph{enhanced} variant adds Double
  DQN~\cite{VanHasselt2016}, a dueling head~\cite{Wang2016}, prioritized experience
  replay~\cite{Schaul2016}, soft target updates, and action masking. These agents cannot handle a
  variable path set and stand in for the original design~\cite{Bourak2025}.
  \item[Path-scoring DQN (simple / enhanced).] The permutation-equivariant design of
  \cref{sec:p1design:scoring}: a shared per-path scorer producing one $Q$-value per path, with variable
  $N$ handled by padding to the batch maximum and masking. The enhanced variant uses the dueling scoring
  head, Double DQN, and prioritized replay. This family delivers the dynamic-action property.
  \item[Conditional path-scoring DQN.] The recommended agent, adding the weight-FiLM dueling head of
  \cref{sec:p1design:film}. Its global state is extended with the five-dimensional encoded intent vector,
  and the FiLM sub-network modulates the per-path features feeding the advantage stream. Two intent
  encodings are supported and recorded in the checkpoint -- the raw weights, or a policy-friendly rescaling
  used in all reported runs.
  \item[Conditional concatenation variants.] Two controls for the FiLM head, differing only in where the
  intent enters. \emph{Value-Concat} appends the intent to the value stream alone, leaving the advantage
  stream on raw path features; by \cref{prop:p1cancel} it cannot re-rank. \emph{Two-Stream-Concat}
  appends the intent to \emph{both} streams, so the advantage stream sees it alongside each path's
  features and the proposition does not apply.
\end{description}

\gap{this list describes the agents the codebase implements, not the agents the reported run actually
trained. The run behind \cref{sec:sp:eval} produced four checkpoints: the enhanced flat DQN and the three
conditional variants. The simple flat agent and \emph{both} unconditioned path-scoring agents were not
trained. That last omission matters most: the enhanced path-scoring DQN is the architectural control that
separates ``conditioning helps'' from ``per-path scoring helps'', and without it the chapter cannot
attribute the roughly 44\% goodput gap between the flat and conditional agents
(\cref{tab:p1eval:ablation}) to conditioning rather than to the scoring architecture. Train it and add it
to the ladder.}

\paragraph{Training details.} All agents use two hidden layers of width 128, learning rate $10^{-3}$,
discount $\gamma = 0.95$, batch size 32, a replay buffer of $10^4$ transitions, and soft (Polyak) target
updates with $\tau = 0.05$. Prioritized replay uses priority exponent $\alpha = 0.6$ with the
importance-sampling exponent $\beta$ annealed from $0.4$ to $1.0$. Exploration is $\varepsilon$-greedy,
decayed multiplicatively per episode. Invalid actions are masked to a large negative value both when
acting and when forming bootstrap targets, so masked paths never contribute to the loss. Episodes are
24~steps (one simulated day); the flat agent trains for 533~episodes and the conditional agents for 666,
with one gradient step every four environment steps. The conditional agents are trained across
reward profiles on a stratified schedule (\cref{sec:p1design:film}), so that within a single training run
the network sees every intent an equal number of times.

\gap{three facts about this training budget belong in the text, because each of them bounds a claim made
later. (i) It is short: 533 and 666~episodes at one gradient step per four environment steps amount to
roughly 3200 and 4000~optimizer updates in total, and $\varepsilon$ decays only to 0.035, never reaching
its 0.01 floor. (ii) The flat DQN's mean training reward is $-0.030$ and its training loss \emph{rises}
over the run, so the weakest rung of the ablation ladder may be undertrained rather than
architecturally limited -- which would inflate the gap attributed to conditioning in
\cref{sec:p1eval:ablation}. (iii) Every result in this chapter comes from a \emph{single seed}. Report a
training-curve figure from the four saved statistics files to substantiate (i) and (ii), and see the
\gap\ in \cref{sec:p1eval:ablation} on seed variance.}

\subsection{Evaluation Pipeline}
\label{sec:p1impl:pipeline}

The system is driven by a six-stage pipeline, each stage writing artifacts consumed by the next:

\begin{enumerate}
  \item \textbf{Topology generation} -- BRITE produces an AS graph, converted into SCION roles, ISDs, and
  links.
  \item \textbf{Beaconing} -- PCB propagation discovers up/core/down segments and composes candidate paths
  per source--destination pair.
  \item \textbf{Traffic simulation} -- 28~days of hourly demand with per-link aggregation, producing the
  time-varying link states the environment replays.
  \item \textbf{Training} -- all agent variants are trained on the resulting environment; checkpoints and
  per-episode statistics are saved.
  \item \textbf{Evaluation} -- every method (learned and baseline) is run over held-out pairs, hours, and
  intents, logging reward, achieved latency and bandwidth, probing overhead, and selection time.
  \item \textbf{Figures} -- results are rendered in a consistent style for the report.
\end{enumerate}

Running the whole pipeline is a single command; intermediate artifacts (topology, path store, link
states, checkpoints, results tables) live under a timestamped run directory so that any stage can be
re-run or inspected in isolation. Key knobs -- topology size, number of training episodes, the cap on
training pairs, and learning hyperparameters -- are exposed as configuration so that experiments can be
scaled up or down.

\subsection{Baselines and Probing Accounting}
\label{sec:p1impl:baselines}

\paragraph{Baselines.} Six heuristic selectors provide the comparison points. \emph{Shortest-path} takes
the fewest AS hops and \emph{lowest-latency} the smallest measured delay. \emph{Widest-path} takes the
largest measured available bandwidth, and is the only baseline requiring a full bandwidth probe on every
candidate. \emph{ECMP} restricts to the minimum-hop set and pins the flow to one member by hashing its
five-tuple, and the \emph{SCION default} takes the minimum hop count with ties broken by latency. Finally
\emph{random} fixes one pseudo-randomly chosen path per source--destination pair and keeps it for the
whole evaluation, which makes it a static-misconfiguration reference rather than a per-decision coin
flip. Together they span the spectrum from naive to strong single-metric heuristics (widest for
throughput, lowest-latency for delay), so that the learned selectors must beat not just a weak baseline
but the \emph{right} heuristic for each intent.

\gap{no \emph{oracle} is evaluated, yet \cref{sec:p1eval:overhead} and the chapter summary both claim to
reproduce a ``near-oracle-quality'' result. Widest-path is a heuristic, not an upper bound: it maximizes
bandwidth, which is the correct target only under the Throughput intent. The right ceiling is a
per-context $\arg\max$ over the true reward under the intent being scored, which would give a genuine
upper bound \emph{for each intent} and turn \cref{tab:p1eval:ablation} into a normalized optimality gap.
A greedy oracle already exists in the codebase but has never been run, and it normalizes goodput by the
per-topology capacity rather than by the per-hour best candidate the agents use, so the two are not
comparable until that is reconciled. Until this is done, replace ``near-oracle'' with ``matching the
strongest heuristic'' throughout.}

\paragraph{Probing accounting.} A fair comparison must charge each method for the measurements it needs.
The environment therefore records, per decision, which paths a method probed and at what cost -- latency
probes being cheap and full bandwidth probes expensive (\cref{sec:p1design:env}). Heuristics that require
a metric on every path (e.g.\ widest-path) incur the cost of probing all paths, whereas the learned
selector is charged only for the paths it actually inspects. This makes the probing-overhead reduction a
directly measurable quantity rather than an assumed benefit, and is the mechanism by which we test the
high-quality-at-low-probing claim of Bourak~et~al.~\cite{Bourak2025}.

\gap{state which baselines are charged which probe type, since it drives the whole of
\cref{fig:p1eval:probing}: widest-path is charged a full bandwidth probe on all 30~candidates, every
other heuristic a latency probe on all 30, and the learned selector a full probe on the single path it
selects. Note also that ECMP is evaluated with a constant flow identifier, so it pins every flow of a
pair to the same path and exhibits none of the load spreading its name implies -- it is effectively
``hash-pinned shortest path'' here and should be labeled as such.}

\section{Experimental Evaluation and Intent Alignment}
\label{sec:sp:eval}

This section evaluates the selectors on held-out hours. We first fix the
experimental setup and metrics (\cref{sec:p1eval:setup}). We then answer three questions the design
raises, in turn: where must intent enter for conditioning to have any effect
(\cref{sec:p1eval:ablation})? Does the conditioned policy re-rank paths in the \emph{right} direction as
the intent changes (\cref{sec:p1eval:intent})? And does it retain the high-quality-at-low-probing
property of the selector it builds on (\cref{sec:p1eval:overhead})? \Cref{sec:p1eval:ceiling} then turns
to the limit that motivates the next chapter.

\subsection{Experimental Setup}
\label{sec:p1eval:setup}

\paragraph{Environment.} Experiments use one BRITE-generated SCION topology of 90~ASes and 528~inter-AS
links, partitioned into 2~ISDs with 6~core ASes. Beaconing (\cref{sec:p1impl:pipeline}) discovers
1004~core segments and 1127~up/down segments, which compose into candidate paths for 8010~routable
source--destination pairs, each with up to 30~paths. The network is driven by the 28-day hourly demand
model with per-link aggregation (\cref{sec:p1design:env}).

\paragraph{Train/test split.} Training uses the first 14~days; all numbers in this chapter are measured
on the last 14~days (336~held-out hours), a window the agents never trained on. Unless noted, each cell
of a results table aggregates 32~pairs across those 336~hours, i.e.\ 10752~independent decision
contexts.

\gap{the split is \emph{temporal only}, and the chapter should say so plainly rather than describe the
pairs as held out. The 32~evaluation pairs are the first entries of the same 8010-pair pool the agents
train on, and all 32 share the same source AS. This is a materially weaker generalization test than the
current wording implies, and it leaves the \emph{Structural \& Topological Generalization} property of
\cref{sec:problem:properties} entirely unevaluated. The fix is cheap and worth doing: hold out a disjoint
set of source--destination pairs spanning several source ASes, and report the same tables on them.}

\paragraph{Intents.} The conditional agent is trained across the six reward profiles of
\cref{sec:p1design:film} and evaluated on the four \emph{selection-relevant} ones:
throughput-maximizing (\emph{Throughput}), latency-averse (\emph{Low-Latency}), loss-averse
(\emph{Low-Loss}), and a balanced midpoint (\emph{Balanced}); their exact weight vectors are given in
\cref{tab:app:profiles}. The two probe-frugality profiles are excluded from the conditioning studies
because a single-path selector issues one probe per decision regardless of $w_{\text{probe}}$, so the
weight has little leverage over the choice; probing cost is instead the subject of
\cref{sec:p1eval:overhead}.

\gap{this exclusion is stated more strongly in the draft than the environment supports. Probe cost is
charged \emph{per hop} (\cref{sec:p1design:env}), so a probe-frugal intent has a real, if weak, incentive
to prefer short paths, and $w_{\text{probe}}$ is not strictly path-independent. The evidence is already
visible in \cref{fig:p1eval:probing}: the learned selector and the random baseline both inspect exactly
one path, yet cost \SI{135}{\milli\second} and \SI{360}{\milli\second} respectively -- a difference that
can only come from path length. These two profiles consume a third of the training schedule and are then
never evaluated; the codebase's multi-profile comparison is the one place their effect on probe count and
path length can be measured, and it has not been run. Either report it or justify the exclusion in terms
of a measured effect size.}

\paragraph{Metrics.} For each method we report the intent-weighted reward of \cref{eq:p1reward} and the
achieved per-flow goodput and latency. Intent alignment is assessed two ways: by reward under each
intent, and by \emph{behavioral divergence} -- the degree to which the chosen path changes as the intent
vector changes, holding the decision context fixed. Concretely, writing $d$ for the number of distinct
paths a method selects across the $K = 4$ intents in a given context and $N$ for that context's candidate
count, the divergence is
\begin{equation}
\mathrm{div} = \frac{d - 1}{\min(K, N) - 1} \in [0,1],
\label{eq:p1divergence}
\end{equation}
averaged over contexts, so a method that ignores intent scores~0 and one that picks a different path for
every intent scores~1. We also report the \emph{choice entropy}, the Shannon entropy in bits of the
distribution of choices across the four intents, which credits a method for spreading its choices evenly
rather than merely for making more than one. Probing cost is measured by the per-selection probe
accounting of \cref{sec:p1impl:baselines}.

\subsection{Where Must Intent Enter? The Conditioning Ablation}
\label{sec:p1eval:ablation}

\Cref{prop:p1cancel} predicts that conditioning routed only through path-independent terms cannot change
the selected path. We test that prediction, and the design claim built on it, by walking up the agent
ladder of \cref{sec:p1impl:agents}: a Flat DQN with a single baked-in objective; a Value-Concat agent
that receives the intent in its value stream alone; and the FiLM agent that modulates the advantage
stream's path features with it. \Cref{tab:p1eval:ablation} reports the mean reward each agent attains
under each intent, together with its adaptivity and choice entropy.

\begin{table}[t]
  \centering
  \caption{Intent-conditioning ablation. Mean intent-weighted reward under each of the four intents;
  \emph{Adaptivity} is the mean behavioral divergence of \cref{eq:p1divergence} (0~=~intent-invariant
  path choice, 1~=~a different path per intent) and \emph{Entropy} the mean choice entropy in bits.
  \emph{Params} counts the $Q$-network's parameters in thousands. Each reward cell aggregates
  10752~decisions over 336~held-out hours, from a single training seed.}
  \label{tab:p1eval:ablation}
  \small
  \begin{tabular}{lccccccc}
    \toprule
    & \multicolumn{4}{c}{Mean reward by intent} & \multicolumn{2}{c}{Intent sensitivity} & \\
    \cmidrule(lr){2-5}\cmidrule(lr){6-7}
    Method & Thpt. & Low-Lat. & Low-Loss & Bal. & Adapt. & Entropy & Params \\
    \midrule
    Flat DQN      & 0.333 & 0.701 & 0.868 & 0.561 & 0.011 & 0.029 & 85.3\,k \\
    Value-Concat  & 0.979 & 0.708 & 0.937 & 0.894 & 0.050 & 0.130 & 36.9\,k \\
    FiLM          & 0.980 & 0.732 & 0.939 & 0.894 & 0.128 & 0.314 & 38.8\,k \\
    \bottomrule
  \end{tabular}
\end{table}

Three findings stand out. First, the unconditioned Flat DQN collapses onto a single trade-off: it scores
well on Low-Loss (0.868) but poorly on Throughput (0.333) and Balanced (0.561), and its adaptivity of
0.011 confirms it selects essentially the same path regardless of the intent it is asked to serve. It is
also the weakest agent in absolute terms, achieving roughly \SI{5.65}{\giga\bit\per\second} of goodput
against the conditional agents' \SI{8.2}{\giga\bit\per\second}, and its reward variance under the
Throughput intent is more than twenty times theirs. Second, both conditional agents recover strong reward
on every intent, because both receive the intent vector. The decisive difference is \emph{how much they
act on it}: FiLM re-ranks paths far more than value-stream concatenation, with an adaptivity of 0.128
against 0.050 and a choice entropy of 0.31 against 0.13~bits. This is the predicted signature of
\cref{prop:p1cancel}: an intent that enters only the value stream shifts all path scores together and
cancels under the $\arg\max$, so Value-Concat can act on intent only through the indirect effect of
training, even though it technically ``takes intent as input.''

Third, the extra re-ranking that FiLM performs is not idle motion: it converts into reward exactly on the
intent that requires re-ranking. On Throughput, Low-Loss, and Balanced the two conditional agents tie,
because those objectives are served by the same high-bandwidth paths the value estimate already favors.
On Low-Latency -- the intent that must steer \emph{away} from the high-bandwidth default toward short,
lightly loaded paths -- FiLM leads, 0.732 against 0.708, and it is the only agent that visibly trades
goodput for delay, giving up about \SI{900}{\mega\bit\per\second} to cut \SI{4.5}{\milli\second} of
latency. We find that FiLM matches or exceeds value-stream concatenation on every intent while never
falling below it, at a cost of 1934~additional parameters, and that the behavioral divergence exposes the
mechanism gap even where the reward gap is small.

\paragraph{The two-stream control.} The ladder above isolates \emph{where} the intent enters, but it does
not by itself establish that modulation is the only workable mechanism, because Value-Concat is
architecturally incapable of re-ranking rather than merely bad at it. The honest control is
Two-Stream-Concat, which appends the intent to the advantage stream as well, so that
\cref{prop:p1cancel} no longer applies. Measured on the same contexts, it attains an adaptivity of
\textbf{0.149} -- \emph{above} FiLM's 0.128.

\gap{this result is currently absent from the chapter and it must be reported. As it stands the draft
compares FiLM against a baseline that cannot re-rank by construction, claims a
$2.5\times$ adaptivity win, and omits the one baseline that can re-rank and that beats FiLM on precisely
that metric. Three things are needed. (i) Recover Two-Stream-Concat's per-intent rewards -- the
checkpoint and its training statistics survive in the run directory, but the reward matrix was
overwritten with the three-method version, so the ablation must be re-run with it included. (ii) Add it
as a fourth row of \cref{tab:p1eval:ablation}. (iii) Rewrite the argument this chapter, the abstract, and
\cref{sec:intro:contributions} make: the defensible claim is not that concatenation fails, but that
conditioning entering only path-independent terms provably cannot re-rank, and that FiLM obtains
comparable adaptivity to per-path concatenation with a structural guarantee and a smaller conditioning
surface. If the two variants are genuinely tied on quality, say so and argue the case on parsimony and
on \cref{prop:p1cancel} -- that is a defensible position, and it is far stronger than an unreported
negative.}

\gap{no statistical claim in this table is currently supported. Every number comes from a single training
seed, and the reported spread is across selections rather than across runs, so it says nothing about
whether re-training would reproduce the ordering. The load-bearing comparison of the whole chapter is a
0.024 reward gap on Low-Latency (0.732 against 0.708). Train each agent under at least five seeds, report
means with confidence intervals, and add a paired significance test over the 10752~shared contexts. The
same applies to the 0.128-against-0.149 adaptivity comparison above, which is well within the range a
seed change could plausibly move.}

\subsection{Intent Alignment}
\label{sec:p1eval:intent}

Adaptivity establishes that the FiLM policy changes its choice with intent; this section shows that it
changes it in the \emph{right} direction. We evaluate the single trained FiLM checkpoint under each
intent and measure both the reward it earns and the metrics of the paths it selects.

\begin{figure}[t]
  \centering
  \includegraphics[width=0.72\linewidth]{p1eval_intent_heatmap.png}
  \caption{Reward matrix $R(\text{intent told}, \text{objective scored})$ for the FiLM policy. Cell
  color is the reward relative to its column mean, so a red diagonal means that conditioning on an intent
  is the best way to score well under that same objective. The diagonal dominates the two objectives that
  trade off against the high-bandwidth default (Throughput, Low-Latency); on Low-Loss and Balanced the
  four intents score alike, because those objectives do not conflict with the default.}
  \label{fig:p1eval:heatmap}
\end{figure}

\Cref{fig:p1eval:heatmap} reports the reward $R(i, j)$ obtained when the agent is told intent~$i$ but its
choice is scored under objective~$j$; each cell is colored by its deviation from the column mean. Reading
down a column asks: to do well under this objective, which intent should the agent be given? For
Throughput the answer is unambiguous -- being told ``Throughput'' scores 0.98, far above any mismatched
intent, and telling the agent the wrong intent (Low-Latency) drops it to 0.75. For Low-Latency the
matching intent again wins its column (0.73), and the mismatch penalty is largest here, precisely because
latency is the objective that most demands re-ranking.

The remaining two columns behave differently, and we report the result rather than the expectation. Under
Low-Loss all four intents lie within 0.01 of one another (0.93--0.94), so the diagonal does not win in any
meaningful sense; under Balanced the diagonal is actually \emph{beaten}, with the Throughput intent
scoring 0.90 against the matched intent's 0.89. The reason is visible in the profiles: low-loss and
low-delay paths in this environment largely coincide with the high-bandwidth paths that most intents
already prefer, so the Low-Loss and Balanced objectives do not conflict with the default and there is
nothing for conditioning to do. Conditioning helps, in other words, exactly where the objectives
genuinely conflict -- which here is two of the four intents we chose.

\gap{this is a weakness in the \emph{choice of intents}, not in the mechanism, and the chapter would be
considerably stronger for fixing it rather than explaining it. The Low-Loss and Balanced profiles produce
near-identical selections on every metric, so the four-intent evaluation is effectively a three-intent
one, and \cref{sec:problem:properties}'s Intent Alignment property is demonstrated on half the intents
reported. Replace one of them with a profile whose optimum genuinely conflicts with high bandwidth -- a
hop-count-averse or trust-dominant intent -- so that all four columns have something to discriminate.
Report the raw reward matrix as a table alongside the figure as well: the column-mean-relative coloring
communicates the pattern but hides that three of the four columns span less than 0.05.}

\begin{figure}[t]
  \centering
  \includegraphics[width=\linewidth]{p1eval_intent_boxplots.png}
  \caption{Distributions of the chosen path's metrics as the FiLM policy is conditioned on each intent,
  over a fixed set of decision contexts. The Low-Latency intent lowers the chosen latency and the
  Throughput intent raises the chosen goodput. Trust separates by intent only weakly, and in a direction
  the trust definition explains rather than the intent labels do.}
  \label{fig:p1eval:boxplots}
\end{figure}

\Cref{fig:p1eval:boxplots} shows the same behavior in the physical metrics of the chosen paths. Holding
the decision contexts fixed and varying only the intent, the Low-Latency intent shifts the chosen-latency
distribution down markedly (median \SI{12.4}{\milli\second} against roughly \SI{16.8}{\milli\second} for
the other three intents), and pays for it in goodput -- \SI{7.31}{\giga\bit\per\second} on average
against \SI{8.15}{\giga\bit\per\second}. The Throughput intent shifts the chosen-goodput distribution up,
though only slightly, since three of the four intents already converge on the widest paths.

The trust panel does \emph{not} behave as the intent labels would suggest, and the discrepancy is
instructive. The highest-trust selections come from the \emph{Low-Latency} intent (mean 0.935), not the
Low-Loss intent (0.919, indistinguishable from Balanced's 0.919 and barely above Throughput's 0.913).
This follows directly from \cref{eq:p1trust}: trust is scored with canonical weights that penalize loss
and delay equally, so the short paths a latency-averse intent selects maximize it, whereas the Low-Loss
profile weights loss heavily and delay barely and therefore optimizes a quantity the canonical trust
score largely ignores. The mechanism is working; the metric and the profile are misaligned.

\gap{the third panel of \cref{fig:p1eval:boxplots} is titled ``Low-Loss intent $\rightarrow$ higher'' and
bolds the Low-Loss box, but the data show Low-Latency winning. Regenerate the figure so the annotation
matches the result -- either by scoring trust under each intent's \emph{own} $w_3, w_4$ rather than
canonical weights, which is the comparison the panel is trying to make, or by replacing the panel with
the chosen path's loss rate, which is what the Low-Loss intent actually optimizes. As it stands the
figure contradicts its own caption.}

\subsection{Zero-Shot Generalization to Unseen Intents}
\label{sec:p1eval:zeroshot}

\gap{this subsection is missing and is the chapter's most conspicuous hole. \Cref{sec:intro:contributions}
claims a policy that serves ``a spectrum of application intents zero-shot'', \cref{sec:problem:properties}
lists \emph{Zero-Shot Generalization Across Objectives} as a required property, and
\cref{sec:analysis:generalization} asserts the FiLM head interpolates between trained profiles. Nothing in
the evaluation tests any of this: all four intents reported above are members of the six-profile training
set, so the results establish only that one network can store six behaviors, not that it generalizes over
the intent space. Two experiments would close it, and both reuse the existing evaluation harness with a
different weight vector. (i) \emph{Interpolation}: sweep $\mathbf{w}$ along the segment between two
trained profiles that genuinely conflict -- Throughput and Low-Latency -- and plot chosen latency and
chosen goodput against the interpolation coefficient. The property holds if the curve is monotone and
passes through the two endpoint behaviors; a step function would show the network has memorized six
points rather than learned a mapping. (ii) \emph{Extrapolation}: evaluate a profile held out of training
entirely, and compare its reward against an agent trained on that profile alone. Report both against the
Value-Concat control, since a conditioning mechanism that fails to interpolate is worth knowing about
regardless of which one is used.}

\subsection{Variable Path Sets}
\label{sec:p1eval:dynamic}

\gap{this subsection is missing. The permutation-equivariant scoring architecture is
\cref{sec:intro:contributions}'s second contribution and the stated reason for abandoning the fixed-action
DQN (\cref{sec:p1design:strawman}), and \cref{sec:problem:properties} lists \emph{Dynamic-Action Handling}
as a required property -- but no result in this chapter is broken down by candidate-path count, so the
claim rests entirely on the architecture's construction. This is the cheapest gap in the chapter to close,
because it is a re-analysis of results already on disk rather than a new run: bin the held-out decisions
by their candidate count $N$ and report selection quality per bin, for the scoring agents and for the
flat agent that must pad to 30. Note the limitation honestly when doing so -- in the reported topology
7988 of 8010~pairs have exactly 30~paths, so the current run has almost no variation in $N$ to measure.
Demonstrating this property convincingly needs either a topology whose pairs differ in path count, or an
evaluation that subsamples the candidate set to a range of sizes at decision time. The second is
straightforward and would also test order invariance directly, by re-scoring the same context under
permuted path orderings and confirming the choice is unchanged.}

\subsection{Probing Overhead}
\label{sec:p1eval:overhead}

A path selector is only useful if its selection quality does not come at the price of measuring every
path on every decision. Using the probing accounting of \cref{sec:p1impl:baselines}, we compare the FiLM
selector against the six heuristics under the Balanced intent so that goodput and reward are directly
comparable.

\begin{figure}[t]
  \centering
  \includegraphics[width=0.72\linewidth]{p1eval_probing.png}
  \caption{Achieved goodput against probing cost per selection (top-left is better). The FiLM selector
  reaches the goodput of exhaustive widest-path selection at roughly \(40\times\) less probing, because it
  inspects the one path it selects rather than all~30. It is also the cheapest method on the plot: every
  load-blind heuristic must measure all 30~candidates to apply its rule, and so probes about
  \(2.7\times\) more than the learned selector while choosing worse.}
  \label{fig:p1eval:probing}
\end{figure}

\Cref{fig:p1eval:probing} places each method in the quality--cost plane. Widest-path attains the highest
goodput (\SI{8.24}{\giga\bit\per\second}) but must issue a full bandwidth probe on all 30~candidate
paths, at \SI{5388}{\milli\second} of probing per selection. The FiLM selector reaches
\SI{8.15}{\giga\bit\per\second} -- within about 1\% of widest-path -- while inspecting only the single
path it selects, at \SI{135}{\milli\second} per selection: a reduction in probing cost of roughly
$40\times$ at matched goodput, and one probe per decision rather than thirty.

The learned selector is in fact the cheapest method in the comparison, not merely cheaper than the
strongest one. Every load-blind heuristic needs its metric on all 30~candidates before it can apply its
rule, so each pays \SI{360}{\milli\second} of latency probing per selection -- about $2.7\times$ the
learned selector's cost -- and then chooses worse: ECMP, the strongest of them, reaches only
\SI{6.6}{\giga\bit\per\second}, and the SCION default, shortest-path, and lowest-latency deliver
\SIrange{5.8}{6.1}{\giga\bit\per\second}, roughly 20--40\% less goodput than FiLM, because none of them
inspects bandwidth. On intent-weighted reward the FiLM selector edges out widest-path, 0.894 against
0.891, despite its slightly lower goodput; the margin comes from the trust term rather than from the
probing penalty, which \cref{sec:p1design:env}'s per-probe normalization keeps comparable between the two.
For the conditioned selector we therefore recover the high-quality-at-low-probing property of
Bourak~et~al.~\cite{Bourak2025}.

\gap{this comparison is run only under the Balanced intent, which is the one intent for which the
selector's advantage is smallest and for which a single heuristic (widest-path) is close to the right
rule. Repeat it per intent: the interesting claim is that one conditioned policy tracks the
\emph{per-intent strongest} heuristic -- widest-path under Throughput, lowest-latency under Low-Latency --
at a fraction of the probing cost of any of them, which is a strictly stronger statement than the one made
here and is what \cref{sec:disc:interpretation} already asserts. Report probes per selection alongside
milliseconds in the figure, since ``1 versus 30'' is the more legible quantity.}

\section{The Single-Path Performance Ceiling}
\label{sec:p1eval:ceiling}

The results so far show that one conditioned policy can serve the evaluated intents at low probing
cost. This final experiment asks the question that motivates the next chapter: how far does the best
single-path choice carry an application as the network fills up? We bin every held-out decision into six
equal-count bins by its realized congestion -- the mean utilization across the flow's candidate paths at
that hour -- and plot the goodput each method achieves against it (\cref{fig:p1eval:ceiling}). Each bin
holds 1792~decisions per method. Because the environment allows links to be oversubscribed, capping
utilization at 1.5 rather than 1.0 (\cref{sec:p1design:env}), the upper three bins sit above 1.0: they are
hours in which the aggregate demand offered to a flow's candidate paths exceeds their capacity.

\begin{figure}[t]
  \centering
  \includegraphics[width=0.72\linewidth]{p1eval_ceiling.png}
  \caption{Achieved goodput against realized congestion (mean path utilization). The FiLM selector tracks
  the widest-path ceiling at every load, but that ceiling itself falls by about 25\% from light to heavy
  congestion: past a point, no single path can carry the demand. Load-blind heuristics degrade far more
  steeply.}
  \label{fig:p1eval:ceiling}
\end{figure}

Two effects appear. First, the FiLM selector tracks the widest-path ceiling across the entire congestion
range, confirming that it makes near-optimal use of whatever single-path capacity exists, while the
load-blind heuristics fall away sharply -- shortest-path, lowest-latency, and the SCION default lose
roughly 55--58\% of their goodput from the lightest to the heaviest congestion bin, ECMP degrades more
gracefully but still loses about 39\%, and random collapses entirely. Second, and crucially, \emph{the
ceiling itself descends}: even FiLM and widest-path drop from about \SI{9.1}{\giga\bit\per\second} at
light load to about \SI{6.8}{\giga\bit\per\second} at heavy load, a loss of roughly 25\%. This is the
single-path performance ceiling. When the shared bottleneck is congested, the best available path is
still one path, and one path can carry only so much; no amount of intent-aware selection recovers the
capacity that a single link cannot provide. Escaping this ceiling requires spreading a flow across
several paths at once -- the problem \cref{ch:multipath} takes up.

\paragraph{Why the reward does not see the ceiling.} The same experiment records each method's
intent-weighted reward, and it tells a complementary story worth stating explicitly: the FiLM selector's
\emph{reward} is essentially flat across the congestion range (0.896 to 0.886) while its goodput falls by
a quarter. This is not a contradiction but a direct consequence of \cref{eq:p1reward}, in which goodput is
normalized against the best candidate path at the same hour. A policy that keeps choosing the best
available path scores near the top whatever that path can carry, so the reward measures selection quality
and is blind to the ceiling by construction. The heuristics, which do not keep choosing the best available
path, collapse on reward as well -- shortest-path from 0.72 to 0.31, and the fixed-random baseline from
0.26 to $-0.32$. The ceiling is therefore invisible to the objective the agent optimizes, which is
precisely why it takes a change of \emph{action space}, rather than a better policy or a better reward, to
break it.

\gap{plot this. \Cref{fig:p1eval:ceiling} shows only the goodput panel, but the reward-versus-congestion
column is already in the same results file and makes the chapter's bridging argument far more forcefully
than the prose above: a two-panel figure in which reward stays flat while goodput descends is the single
clearest statement of ``the policy is doing its job and the architecture is the limit''. Add reward as a
second panel and give both panels the bin-count annotation.}

\paragraph{Summary.} This chapter established four claims. One conditioned scoring policy serves the four
evaluated intents with a single model, matching on each the heuristic that is strongest for it. Making
that conditioning effective requires routing it to the per-path features: by \cref{prop:p1cancel}, intent
that enters only path-independent terms cancels under the $\arg\max$, and the value-stream variant duly
re-ranks $2.5\times$ less than FiLM while losing on the intent that most demands re-ranking. The learned
selector reaches widest-path goodput while probing roughly $40\times$ less -- and less than every other
baseline in absolute terms -- recovering the high-quality-at-low-probing property it builds on. And even
this strong single-path selector loses about a quarter of its goodput under heavy congestion, because a
single path is a hard capacity bound -- the ceiling that \cref{ch:multipath} sets out to break.

\gap{two of these four claims are narrower than the chapter's framing suggests, and the framing should be
brought down to them until the corresponding experiments exist. ``Serves the intent spectrum'' is
currently ``serves four trained profiles, two of which are behaviorally indistinguishable'', pending
\cref{sec:p1eval:zeroshot}. ``Requires modulation'' is currently ``requires not routing intent through
path-independent terms'', pending the two-stream control of \cref{sec:p1eval:ablation}. Both are still
worthwhile results; they are just not the results the abstract and \cref{sec:intro:contributions} promise.}
